            Parameter {\bf in} is the queue-mode input variable.
            
            The queue must be written somewhere in the program (queues without
            sources result in a fatal error on completion of immediate
            execution). When the module is created the argument queue must have
            an established write clock domain or a fatal error is generated.

            If the input queue is read anywhere else in the program (i.e. other
            than in this module) then it must be read in the bus clock domain
            since that is the clock domain in which it is read in the module
            created by this procedure.

            If optional parameter {\bf wordcount} has an argument and it
            is true, the number of words in the queue buffer can be interrogated
            by the CPU. The input queue may have any of the allowed buffer sizes.
            If the queue is asynchronous, i.e. the read domain is not the bus
            clock domain, the count will not be exact but will be a guaranteed
            minimum number of readable data.

            If optional parameter {\bf wordcount} has an argument and it
            is false, the data returned from reading the number of words in the
            buffer is 0 or 1 instead of a count. A return of 1 means the queue
            buffer is not empty, i.e. can be read. This option is provided to
            simplify the logic where a simple empty/non-empty status is required.

            If optional parameter {\bf wordcount} has no argument then
            there will no interface to directly read the number of words in the
            queue buffer, however an interrupt may be used instead.            

            Note that an empty queue is guaranteed to return all ones to the CPU
            when read. For most data types this single value will usually be
            invalid and hence in that case it can be used as an indication of an
            empty queue. This can be used instead of explicitly reading the
            available word count.

            If an argument to optional parameter {\bf ilevel} is provided then a
            level interrupt on queue buffer status is established as follows -

            If there is no argument then no interrupt is provided.

            If the argument is type "log" with value false then the interrupt will
            be asserted when the queue buffer is not empty.

            If the argument is type "log" with value true then the interrupt will be
            asserted when the number of words in the buffer equals or exceeds a
            level which has been written to the interface from the CPU. The initial
            value of this level is 1. A level value of 0 should not be written to
            the interface as that will result in a continuous level interrupt
            regardless of queue state.

            If the argument is type "uint" or "int" then the interrupt will be
            asserted when the number of words in the buffer equals or exceeds the
            fixed value specified by the integer argument.

            The optional output parameter {\bf count} provides a count of the
            number of data in the queue. The inbuilt function queuewords() can
            only be used in a module which reads the queue, which is usually not
            the case where {\bf cpu\_api.read\_queue()} is the queue destination.
            {\bf count} must have clock domain {\bf bus\_clk}. If it has no clock
            domain, {\bf cpu\_api.read\_queue()} will set it to {\bf bus\_clk}.
            The width of {\bf count} must be sufficient to hold the input queue
            depth.

            If optional output parameter {\bf out} has an argument, each datum
            read from the input queue {\bf in} to the CPU local bus is also
            written to queue {\bf out}. The type of {\bf out} must be such that
            it can be assigned from {\bf in}. {\bf out} should not be allowed to
            block since it cannot stop data being popped from {\bf in} to the CPU
            bus. While {\bf out} in unavailable for writing, no attempt is made
            to write data popped from {\bf in} to it.

            If optional output parameter {\bf access} has an argument, the static
            variable of type "log" is assigned true for one clock cycle every
            time a datum is popped from the input queue {\bf in}. {\bf access}
            may have any clock domain.
            
            On running info2c on the .json file the following C functions will have
            been defined for interface {\em name} -
            
            \begin{tabular}{| l | l | l |}
            \hline
            function                                 & action when?           &                         \\ \hline \hline
            fpga\_{\em name}\_read()                 & always                 & read a datum            \\ \hline
            fpga\_{\em name}\_avail\_read()          & wordcount has argument & read number of words    \\ \hline
            fpga\_{\em name}\_avail\_thresh\_write() & ilevel "log" true      & set interrupt threshold \\ \hline
            fpga\_{\em name}\_enable()               & ilevel has argument    & enable interrupt        \\ \hline
            fpga\_{\em name}\_disable()              & ilevel has argument    & disable interrupt       \\ \hline 
            fpga\_{\em name}\_wait()                 & ilevel has argument    & wait on interrupt       \\ \hline  
            \end{tabular}
            
            Example: \\
            If 3PL source file prog.3pl contains the following \\
            \begin{lstlisting}[gobble=12, language=threepl]
                class(cpu, cpu_api());
                
                queue(p, "uint:16");
                .
                .
                cpu.read("P1", p, wordcount=false, ilevel=false);
                .
                .
            \end{lstlisting}
            the CPU code in C should contain something like \\
            \begin{lstlisting}[gobble=12, language=C]
                #include "prog_fpga.h"
                .
                .
                fpga_t *fpga = NULL;
                int     d[...];
                int     i = 0;
                .
                .
                fpga_open();
                .
                .
                fpga_P1_event_wait(fpga);
                while (fpga_P1_avail_read(fpga) != 0)
                    d[i++] = fpga_P1_read(fpga);
                .
            \end{lstlisting}
            The code will wait until the queue p goes from empty to non-empty
            when an interrupt will release the wait. A loop while the queue
            remains non-empty reads until the queue is empty. Note that since
            {\bf wordcount} is given an argument but the value is false, the count
            will only have values 0/1 for empty/not-empty.
           
            This procedure creates no in-line target code so the current clock
            domain when this procedure is called is not relevant. The clock domain
            is not changed by this procedure,
            
            The CPU code in Python3 should contain something like \\
            \begin{lstlisting}[gobble=12, language=python]
                import sys
                sys.path.append("/Users/......./install/3pl")
                import info2py
                from cpu_serial import *

                fpga_open()

                fpga = info2py.info_from_file("prog")   # load info from prog.info file

                d = []
                fpga.P1_wait()
                while (fpga.P1_avail_read() != 0)
                    d.append = fpga.P1_read()

            \end{lstlisting}
            
            Note that in the unusual case where a queue variable is both read and
            written by the CPU and is not otherwise accessed in the FPGA the input
            and output clock domains of the queue will be undefined. To overcome
            this the clock domain can be set when the queu is created, e.g.
            
            \begin{lstlisting}[gobble=12, language=threepl]
            queue(q, "uint10", domain=null);
            \end{lstlisting}
            
            which will set both the input (write) and output (read) clock domains
            to the current clock domain - see \autorefph{sec:ipqueue}).
